\subsection{Cuerpos de definición.}
\label{subsection:IV.4.8-fr}

\begin{env}[4.8.1]
\label{IV.2.4.8.1-fr}
Dado un preesquema $X$, designaremos, en este número, por
$\mathfrak S(X)$ el conjunto de los subpreesquemas de $X$, y, como de costumbre,
por $\mathfrak P(X)$ el conjunto de los subconjuntos del espacio subyacente a $X$;
para todo $\mathcal O_X$-módulo cuasicoherente $\mathscr F$, designaremos
por $\Phi(\mathscr F)$ el conjunto de los $\mathcal O_X$-submódulos
cuasicoherentes de $\mathscr F$.
\end{env}

\begin{env}[4.8.2]
\label{IV.2.4.8.2-fr}
Sean $k$ un cuerpo, $X$, $Y$ dos $k$-preesquemas, $\mathscr F$, $\mathscr G$
dos $\mathcal O_X$-módulos cuasicoherentes, $K$, $K'$ dos extensiones de $k$
tales que $K'\subset K$. Se tienen entonces, correspondientes al morfismo
$\operatorname{Spec}(K)\to\operatorname{Spec}(K')$, las aplicaciones
canónicas
\begin{align}
\Phi(\mathscr F\otimes_kK')
&\longrightarrow\Phi(\mathscr F\otimes_kK)
\tag{4.8.2.1}
\label{IV.2.4.8.2.1-fr}\\
\operatorname{Hom}(\mathscr F\otimes_kK',\mathscr G\otimes_kK')
&\longrightarrow
\operatorname{Hom}(\mathscr F\otimes_kK,\mathscr G\otimes_kK)
\tag{4.8.2.2}
\label{IV.2.4.8.2.2-fr}\\
\mathfrak S(X_{(K')})&\longrightarrow\mathfrak S(X_{(K)})
\tag{4.8.2.3}
\label{IV.2.4.8.2.3-fr}\\
\operatorname{Hom}_{K'}(X_{(K')},Y_{(K')})
&\longrightarrow\operatorname{Hom}_K(X_{(K)},Y_{(K)})
\tag{4.8.2.4}
\label{IV.2.4.8.2.4-fr}\\
\mathfrak P(X_{(K')})&\longrightarrow\mathfrak P(X_{(K)})
\tag{4.8.2.5}
\label{IV.2.4.8.2.5-fr}
\end{align}
Si $p_X:X_{(K)}\to X_{(K')}$ es la proyección canónica, la aplicación
(4.8.2.5) es $M\leadsto p_X^{-1}(M)$, y asimismo (4.8.2.3) es
$Z\leadsto p_X^{-1}(Z)=Z\otimes_{K'}K$ (\textbf I, 4.4.1). La aplicación (4.8.2.1)
es
$\mathscr H\leadsto p_X^*(\mathscr H)
=\mathscr H\otimes_{\mathcal O_{X_{(K')}}}\mathcal O_{X_{(K)}}$;
(4.8.2.2) es $u\leadsto p_X^*(u)$ y, finalmente, (4.8.2.4) es la aplicación
$f\leadsto f_{(K)}$.

Si se denota por $\varepsilon_{K,K'}$ cualquiera de estas aplicaciones
canónicas, es claro que, si $K''$, $K'$, $K$ son tres extensiones de $k$
tales que $K''\subset K'\subset K$, se tiene
\begin{equation}
\varepsilon_{K,K''}=\varepsilon_{K,K'}\circ\varepsilon_{K',K''}.
\tag{4.8.2.6}
\label{IV.2.4.8.2.6-fr}
\end{equation}
\end{env}

\oldpage[IV-2]{81}

\begin{proposition}[4.8.3]
\label{IV.2.4.8.3-fr}
Las aplicaciones canónicas \textup{(4.8.2.1)} a \textup{(4.8.2.5)} son inyectivas.
\end{proposition}

\begin{proof}
En efecto, se sabe que $p_X$ es \emph{fielmente plano} y \emph{cuasicompacto}; el hecho de que
(4.8.2.5) sea inyectiva resulta, por tanto, de que $p_X$ es sobreyectivo.
La inyectividad de (4.8.2.1) resulta de (2.2.2), la de (4.8.2.2) de (2.2.7),
la de (4.8.2.3) de (2.2.15) y, finalmente, la de (4.8.2.4) de (2.2.16).
\end{proof}

\begin{definition}[4.8.4]
\label{IV.2.4.8.4-fr}
Con las notaciones de \textup{(4.8.2)}, se dice que un
$\mathcal O_{X_{(K)}}$-submódulo cuasicoherente de
$\mathscr F\otimes_kK$ \textup{(resp.} un homomorfismo
$\mathscr F\otimes_kK\to\mathscr G\otimes_kK$, \textup{resp.} un subpreesquema de
$X_{(K)}$, \textup{resp.} un $K$-morfismo $X_{(K)}\to Y_{(K)}$, \textup{resp.} un
subconjunto de $X_{(K)}$\textup{)} está \emph{definido sobre $K'$} si pertenece a
la imagen de la aplicación \textup{(4.8.2.1) (resp. (4.8.2.2), resp. (4.8.2.3), resp.
(4.8.2.4), resp. (4.8.2.5))}. Se dice entonces que $K'$ es un cuerpo de
definición del objeto considerado.
\end{definition}

Es claro que el propio $K$ es un cuerpo de definición de un elemento de
cualquiera de los cinco conjuntos de llegada de las aplicaciones (4.8.2.1) a
(4.8.2.5). Se observará, sin embargo, que, para un $K$-preesquema $Z$, decir,
por ejemplo, que un $\mathcal O_Z$-módulo cuasicoherente $\mathscr H$ está
«~definido sobre un subcuerpo $K'$ de $K$~» solo tiene sentido si $Z$
\emph{ha sido dado} en la forma $X_{(K)}$ para un preesquema $X$ sobre un
subcuerpo $k$ de $K$, y si $\mathscr H$ es un $\mathcal O_Z$-submódulo de un
$\mathcal O_Z$-módulo que \emph{ha sido dado} en la forma
$\mathscr F\otimes_kK$, donde $\mathscr F$ es un $\mathcal O_X$-módulo
cuasicoherente; se hacen observaciones análogas para las demás nociones. De
(4.8.2.6) se sigue que, si un elemento de uno de los conjuntos de llegada de
las aplicaciones de (4.8.2) está \emph{definido sobre} $K'$, también está definido
sobre todo cuerpo $K''$ tal que $K'\subset K''\subset K$. Finalmente, con las
notaciones de (4.8.2.6), si $K_1$ es una extensión de $K$, para que un elemento
del conjunto de llegada de $\varepsilon_{K,K'}$ esté definido sobre $K'$, es
necesario y suficiente que su imagen por $\varepsilon_{K_1,K}$ esté definida
sobre $K'$, en virtud de la relación
$\varepsilon_{K_1,K'}=\varepsilon_{K_1,K}\circ\varepsilon_{K,K'}$.
